% ============================================================
%  Chapter 1 — Introduction
% ============================================================
\chapter{Introduction}
\label{chap:intro}

\section{Context and Motivation}

Electricity markets have undergone a fundamental transformation over the past three decades. The progressive liberalisation of European energy markets, accelerated by successive EU directives from 1996 onwards, replaced vertically integrated state monopolies with competitive wholesale markets in which prices are determined by the interaction of supply and demand. In France, the day-ahead electricity market operates through EPEX SPOT, where participants submit bids and offers for each hour of the following day, with a single clearing price determined by a batch auction that closes at 12:00 CET.

The resulting price series is among the most complex in quantitative finance. Unlike commodity prices, electricity cannot be economically stored at scale, which forces supply and demand into instantaneous equilibrium at every moment. The consequence is a price process characterised by extreme volatility, sharp intraday seasonality, strong weekly periodicity, and irregular price spikes that can reach several thousand euros per megawatt-hour during scarcity episodes. The French market adds a layer of specificity: with approximately 63 GW of installed nuclear capacity representing roughly 70\% of annual electricity production \citep{rte2023}, the availability of the nuclear fleet is a primary determinant of price levels. In parallel, France's residential heating stock is predominantly electric, creating a thermosensitivity of approximately 2.4 GW per degree Celsius in winter \citep{rte2024} — the highest in Europe.

Accurate forecasting of day-ahead prices is therefore not merely an academic exercise. It is a critical input for energy suppliers setting customer tariffs, for industrial consumers optimising consumption schedules, for renewables producers valuing flexibility options, and for traders positioning in forward markets. A one euro per megawatt-hour improvement in forecast accuracy over a 100 MW portfolio translates to approximately 876,000 EUR in annual value capture — motivating sustained interest in more accurate models.

\section{Research Question}

This thesis addresses the following central question:

\begin{quote}
\textit{Do weather variables significantly improve the accuracy of day-ahead
electricity price forecasts for the French market, and does their contribution
depend on the market regime?}
\end{quote}

This question is decomposed into four sub-questions:

\begin{enumerate}
  \item Can machine learning models (Random Forest, XGBoost) substantially
  outperform a na\"{i}ve statistical benchmark on French day-ahead prices?
  \item Does the addition of meteorological features (temperature, wind speed,
  solar radiation, Heating Degree Days, Weather Stress Index) produce a
  statistically significant improvement over a model built on market
  fundamentals alone — and is this result stable across market regimes?
  \item Does the marginal value of weather features differ between a stable
  market regime (2024--25) and an extreme-stress regime (2022 energy crisis),
  and what structural mechanisms explain any difference?
  \item Can the predictive improvement of these models be translated into
  economically meaningful returns through a day-ahead directional trading
  strategy with realistic transaction costs?
\end{enumerate}

\section{Contribution}

This thesis makes four contributions to the electricity price forecasting
literature:

\begin{itemize}
  \item \textbf{Comprehensive French dataset:} We assemble and publish a
  reproducible dataset spanning 2018--2025 that combines ENTSO-E generation
  and load data, ERA5 reanalysis weather data, cross-border flow data, and
  fuel price series. To our knowledge, this is one of the most complete
  publicly-sourced datasets for the French day-ahead market.

  \item \textbf{Rigorous weather ablation study with regime comparison:}
  We design a four-model experiment (A/B/C/D) isolating the marginal
  contribution of weather features using a Diebold-Mariano test with
  Harvey-Leybourne-Newbold correction, applied to two out-of-sample periods.
  Our key finding is that weather is statistically non-significant in a stable
  regime ($p = 0.572$) but highly significant in the 2022 energy crisis
  ($p < 0.001$, DM $= -13.27$). We explain this regime-dependence through
  the \textit{information redundancy hypothesis}: in normal conditions, the
  RTE load forecast and TTF gas prices subsume the weather signal; in crisis
  conditions, the gas supply shock breaks this redundancy.

  \item \textbf{2022 energy crisis as a natural experiment:} We exploit the
  2022 French-European energy crisis — characterised by simultaneously low
  nuclear availability ($\approx$40\%) and extreme gas prices — as a stress
  test of the weather redundancy hypothesis. This provides the first formal
  evaluation of regime-dependence in the weather-price relationship for the
  French market.

  \item \textbf{Day-ahead executable trading backtest:} We implement a
  directional trading strategy using an economically coherent signal design
  grounded in the actual mechanics of the EPEX SPOT day-ahead auction, with
  dead-band filtering and three realistic transaction cost scenarios drawn
  from the 2024 EPEX fee schedule.
\end{itemize}

\section{Structure of the Thesis}

The remainder of this thesis is organised as follows. Chapter~\ref{chap:lit} reviews the literature on electricity price forecasting, covering both statistical and machine learning approaches. Chapter~\ref{chap:data} describes the data sources, feature engineering process, and exploratory analysis. Chapter~\ref{chap:model} presents the modelling framework, results, and statistical tests. Chapter~\ref{chap:backtest} introduces the trading backtest and its results. Chapter~\ref{chap:discussion} discusses the findings, limitations, and extensions. Chapter~\ref{chap:conclusion} concludes.

All code, data pipelines, trained models, and figures are available in the project repository at \url{https://github.com/leogthub/fr-power-thesis}.
